\chapter{Referencial Teórico}  \label{cap:03}

Nessa sessão você deve realizar o seu referencial teórico, que consiste na construção de ideias fundamentadas na área que você está estudando.
É nessa sessão que você irá apresentar os conceitos que estará estudando para poder construir sua revisão bibliográfica ou seu software \cite{howlett_parsing_2022}.

Por exemplo:
Se meu estudo é sobre aplicativos de celular para a educação matemática, posso criar uma seção sobre os aplicativos, sobre seu histórico de uso, como são usados na educação, quais são os resultados promissores e desvantagens do seu uso.

Posteriormente, poderia criar um subtópico falando a respeito de como os aplicativos já criados estão atuando na área de informática, quais são as áreas de aplicação, o que esses aplicativos atendem ou deixam de atender na área em que atuam.

Mais adiante, posso criar um capítulo sobre as tecnologias utilizadas para a concepção do projeto, sendo elas ambientes de programação, ambientes de documentação do sistema e projeções de testes.

Veja bem, todas essas áreas pedem muita leitura sobre o tema, para que seja realizada uma boa fundamentação teórica do assunto.

Não esqueça de referenciar suas imagens no texto para que o Overleaf possa montar o índice de imagens. O código da \autoref{fig:rotuloImagem} abaixo insere imagens. Ele está melhor explicado no tutorial de LaTeX.\\

% comando para inserir figura
\begin{figure}[!h]
    \centering
    \includegraphics[width=0.7\linewidth]{imagens/exemplo.png}
    \caption{Legenda da Imagem}
    \label{fig:rotuloImagem}
\end{figure}
\newpage 

Não esqueça de referenciar suas tabelas no texto para que o Overleaf possa montar o índice de tabelas. O código da \autoref{tab:tabela} abaixo insere uma tabela que você pode gerar utilizando o gerador de tabelas. (\href{https://www.tablesgenerator.com/}{Clique aqui para Acessar o Gerador de Tabelas}). Esse passo a passo está melhor explicado no tutorial de LaTeX.\\

% comando para inserir tabelas
\begin{table}[!h]
\centering
\begin{tabular}{l|l|l}
\hline
\multicolumn{1}{c|}{\textbf{Tabela}} & \multicolumn{1}{c|}{\textbf{Título 1}} & \multicolumn{1}{c}{\textbf{Título 2}} \\ \hline
Assunto 1                            & A                                      & C                                     \\ \hline
Assunto 2                            & B                                      & D                                     \\ \hline
\end{tabular}
\caption{Tabela}\label{tab:tabela}
\end{table}

